%%% PREVIEW MARKER (delete at integration)
\par\medskip\noindent\textbf{[Part A. For Section~\ref{sec:ground-correspondence}, first subsection, after the paragraph that ends ``Neither R4 nor R5 is required for the proposition.'']}\par\smallskip

Under these conditions the output ground is an eigenvector of the output return as well. By (v), $A^*y=ax$, and $x\in\operatorname{Dom}A$, so $y\in\operatorname{Dom}(AA^*)$ and
\begin{equation}
 AA^*y=aAx=a^2y .
 \label{eq:ground-output-return}
\end{equation}
By (i), the two grounds are eigenvectors of $A^*A$ and $AA^*$ with the same positive eigenvalue $a^2$. The weak bridge alone does not give this. On the unequal-cell path below, $AA^*y$ is not a multiple of $y$.

%%% PREVIEW MARKER (delete at integration)
\par\medskip\noindent\textbf{[Part B. For Section~\ref{sec:determination}, first subsection, after the paragraph that ends ``no normalized vector $Xy/\|Xy\|$.'']}\par\smallskip

A bounded operator $X$ on $\mathcal H_\Phi$ maps the ground line into itself exactly when
\begin{equation}
 (I-P_y)XP_y=0,\qquad\text{equivalently}\qquad Xy=\omega(X)\,y ,
 \label{eq:ground-line-kept}
\end{equation}
that is, when the variance in Eq.~\eqref{eq:determination-variance} vanishes. It then keeps the ground ray, in the sense that $Xy/\|Xy\|$ is a phase times $y$, exactly when $\omega(X)\ne0$. The condition is one-sided. For a unit vector $w\perp y$, the operator $\ket{y}\!\bra{w}$ sends $y$ to zero and maps the line into itself, while its adjoint sends $y$ to $w$. Both $X$ and $X^*$ map the line into itself exactly when $X$ commutes with $P_y$, because the two conditions remove the two off-diagonal blocks of $X$ with respect to $\C y\oplus y^\perp$. The line then reduces $X$. For a unitary $U$ one condition gives the other, since $Uy=cy$ implies $U^*y=\bar cy$. A bounded self-adjoint operator maps the line into itself exactly when $y$ is an eigenvector.

%%% PREVIEW MARKER (delete at integration)
\par\medskip\noindent\textbf{[Part C. For Section~\ref{sec:asymmetric-nonduality-moments}, subsection ``Fixing the ground and attracting to it'', after the paragraph that ends ``reduced description of a unitary evolution on a larger space.'']}\par\smallskip

For a strongly continuous unitary group $U_t=e^{-itG}$ with self-adjoint generator $G$, the following conditions are equivalent:
\begin{equation}
 \begin{gathered}
 \text{(i) }U_tP_yU_t^*=P_y\ \text{for every real }t,\qquad
 \text{(ii) }U_ty=e^{-itg}y\ \text{for some real }g\text{ and every }t,\\
 \text{(iii) }y\in\operatorname{Dom}G\ \text{and}\ Gy=gy .
 \end{gathered}
 \label{eq:relaxation-unitary-group}
\end{equation}
If (i) holds, then $U_ty=\chi(t)y$ with $|\chi(t)|=1$. The group law and strong continuity make $\chi$ a continuous homomorphism from $\R$ into the unit circle, so $\chi(t)=e^{-itg}$ for a real $g$. By the spectral theorem the derivative of $U_th$ at $t=0$ exists exactly when $h\in\operatorname{Dom}G$, and it is then $-iGh$. For $h=y$ the derivative is $-igy$, which gives (iii). Conversely, the spectral theorem gives $e^{-itG}y=e^{-itg}y$ for an eigenvector. The times at which the projection returns, $U_tP_yU_t^*=P_y$, are those with $|\langle y,U_ty\rangle|=1$. By the group law they form a subgroup of $\R$, and it is closed because $t\mapsto\langle y,U_ty\rangle$ is continuous. A closed subgroup of $\R$ is $\{0\}$, a lattice $c\mathbb Z$ with $c>0$, or $\R$ itself. A single return is therefore not enough: for $G=\operatorname{diag}(0,1)$ and $y=(1,1)/\sqrt2$ the projection returns exactly at the multiples of $2\pi$. Two return times with an irrational ratio are enough, because they generate a dense subgroup. The equivalence is Eq.~\eqref{eq:graph-evolution-phase} for a single vector, and Section~\ref{sec:symmetry-average} uses the same character for a group that leaves the pure state of a vector invariant.

%%% PREVIEW MARKER (delete at integration)
\par\medskip\noindent\textbf{[Part D. For Section~\ref{sec:asymmetric-nonduality-moments}, subsection ``Relaxation across many moments'', after the paragraph that ends ``The second concerns the particular states $\sigma_n$.'']}\par\smallskip

Each projection $P_{y_n}$ has rank one, since it projects onto a vector. LG asks for more. When $y_n$ is the ground of a Hamiltonian $H_n$ at moment $n$, the projection $P_{y_n}$ is the whole spectral projection of $H_n$ at the bottom of its spectrum, as in Eq.~\eqref{eq:ground-spectral-projection}. This is the uniqueness part of R2 imposed at that moment. For a degenerate bottom, $P_{y_n}$ would be one choice inside a larger spectral projection.

%%% PREVIEW MARKER (delete at integration)
\par\medskip\noindent\textbf{[Part E. For Section~\ref{sec:cohomology-index}, first subsection, after the paragraph that ends ``were specified independently.'']}\par\smallskip

The differential has no second step. Its image lies in the output summand, and it vanishes there:
\begin{equation}
 d_A(x,y)\in0\oplus\mathcal H_\Phi,\qquad d_A(0,y)=0,\qquad d_A^2=0 .
 \label{eq:index-no-second-step}
\end{equation}
The last identity holds on the whole domain $\operatorname{Dom}A\oplus\mathcal H_\Phi$, because $d_A$ maps that domain into itself. The same holds for $-d_A$. In this sense the model generates no sequence $x\mapsto Ax\mapsto A(Ax)\mapsto\cdots$. An output vector is not an input to the same map. The opening section takes a cognition to be manifest in its own occurrence. On that reading no second act is needed to make it manifest. Whether the vanishing of $d_A$ on the output summand expresses this reflexivity is a further identification.

The identity $d_A(0,y)=0$ holds for every $A$ and before any evolution is chosen. It does not describe an evolution that removes $y$. It records that the map is not applied to output vectors. It places no condition on how a channel treats the grounds. In the shift example below, the channels that attain Eq.~\eqref{eq:index-channel-coherence} share one preparation $A$ for every value of $a$. For $a=0$ the channel carries the input ground $f_0$ to the output ground $g_1$. For $a>0$ it sends the input ground to the missing line $g_0$ with probability $a$.

The returns vanish only on the kernels. For $x\in\operatorname{Dom}(A^*A)$,
\begin{equation}
 \langle x,A^*Ax\rangle=\|Ax\|^2,\qquad
 \ker(A^*A)=\ker A,\qquad \ker(AA^*)=\ker A^* .
 \label{eq:return-kernels}
\end{equation}
No bridge is needed for this. It constrains any attempt to take the returns as the sector Hamiltonians. If $H_\Psi=A^*A$ had the ground $x$ with energy zero, then $Ax=0$, contrary to the weak bridge. An eigenvalue of $A^*A$ at $x$ would equal $\langle x,A^*Ax\rangle=a^2$, so $x$ is an eigenvector of $A^*A$ only under strong compatibility, condition (i) of Proposition~\ref{prop:ground-polar-compatibility}. Without it $x$ is a ground of no shift of $A^*A$. LG proposes the shifted return $A^*A-a^2I$ on the active input $S$. Under strong compatibility $x\in S$ and $A^*Ax=a^2x$, so $x$ is a ground of this operator exactly when $a^2$ is the bottom of the spectrum of $A^*A$ on $S$, and the unique ground when that eigenvalue is also simple. The restriction to $S$ matters. For $A=\operatorname{diag}(0,1,2)$ and $x=y=(0,1,0)$ the grounds are strongly compatible with $a=1$. On the whole input space $A^*A-I$ has the eigenvalues $-1,0,3$, and the kernel of $A$ lies below $x$. On $S$ the eigenvalues are $0$ and $3$, and $x$ is the unique ground. On the output side, R3 and R5 already exclude $AA^*$, shifted or not, as a Hamiltonian whose ground is $y$. By R5 the bottom of the spectrum of $AA^*$ is the eigenvalue zero, with eigenspace $\ker A^*\ne\{0\}$. By R3 the vector $y=Ax/a$ lies in the range of $A$ and is orthogonal to that whole eigenspace. Strong compatibility adds that $y$ is an eigenvector, with the eigenvalue $a^2$ of Eq.~\eqref{eq:ground-output-return}, and then $AA^*-a^2I$ equals $-a^2I$ on the defect. In the examples the returns select other states. On the unequal-cell path $A^*A=\operatorname{diag}(1/4,1)$, whose ground is the first coarse site and not the constant ground of the two-site path. In the pair preparation $A_\theta$ is an isometry, so $A^*A=I$ selects no ray. On the two-particle space $A_\theta A_\theta^*$ is the projection onto $D$, with the four-dimensional complement $\C d\oplus W$ as its zero eigenspace.
